% !TEX root = ../notes_template.tex
\chapter{Ankle \& Foot}\label{chp:ankle_foot}

\section*{Objectives}

\begin{itemize}
  \item Analyze ankle–foot control and mobility using top-tier SFMA patterns (ADS, SLS, MSF, MSR).
  \item Perform manual muscle testing (MMT) for dorsiflexion, plantarflexion, inversion, eversion, and toe flexion.
  \item Identify the primary muscle(s), peripheral nerve, and spinal root(s) for each tested action.
  \item Perform and interpret the Achilles tendon reflex to assess tibial nerve and S1 root function.
  \item Identify the role of tibialis anterior and tibialis posterior in dorsiflexion, inversion, and arch stabilization.
  \item Acquire ultrasound images of the Achilles tendon in long- and short-axis at the insertion, mid-substance, and myotendinous junction.
  \item Perform dynamic scanning of the lateral ankle ligaments (ATFL/CFL) during passive inversion and plantarflexion.
  \item Palpate and image the dorsalis pedis or posterior tibial pulse, confirming vascular identity via compressibility and pulsatility.
\end{itemize}

\section*{Functional Integration: Distal Control and Dynamic Support}

The ankle and foot provide dynamic interaction with the ground and are essential for shock absorption, propulsion, and balance. As the distal controller of the lower kinetic chain, this region plays a critical role in load transfer and movement modulation. Dysfunction here may present locally or contribute to proximal issues at the knee, hip, or spine. 

Observation of foot posture, ankle range, and dynamic support during gait or squat activities can inform manual testing priorities. Restoration of ankle dorsiflexion, medial arch integrity, and neuromuscular control of the ankle evertors is often critical for optimizing closed chain movement.

\subsection*{SFMA Top Tier Patterns with Ankle/Foot Relevance}

\begin{itemize}
  \item \textbf{Arms Down Squat (ADS)} \textendash{} Observe for dorsiflexion range, foot pronation/supination, and tibial progression. 
  \item \textbf{Single Leg Stance (SLS)} \textendash{} Assess balance, toe flexor engagement, and control of hindfoot/forefoot alignment. 
  \item \textbf{Multisegmental Flexion (MSF)} \textendash{} Look for restrictions in posterior chain length including gastrocnemius or soleus tightness. 
  \item \textbf{Multisegmental Rotation (MSR)} \textendash{} Monitor foot pivot and forefoot rotation contributions to whole-body turning.
\end{itemize}

\section*{Ankle and Foot Movements: Muscles and Innervation}

\begin{table}[H]
\centering
\captionsetup{justification=raggedright, singlelinecheck=false}
\begin{tabular}{|p{3.5cm}|p{4.5cm}|p{4cm}|p{2.5cm}|}
\hline
\textbf{Movement} & \textbf{Muscle} & \textbf{Peripheral Nerve} & \textbf{Nerve Root(s)} \\
\hline
Plantarflexion & Gastrocnemius, Soleus & Tibial & S1--S2 \\
\hline
Dorsiflexion & Tibialis Anterior & Deep Fibular & L4--L5 \\
\hline
Inversion & Tibialis Posterior & Tibial & L4--L5 \\
\hline
Eversion & Fibularis Longus, Brevis & Superficial Fibular & L5--S1 \\
\hline
Great Toe Extension & Extensor Hallucis Longus & Deep Fibular & L5 \\
\hline
\end{tabular}
\caption{Primary Ankle and Foot Movements with Neuromuscular Innervation}
\end{table}

\begin{tcolorbox}[colback=gray!5, colframe=black, title={Clinical Insight: Tibialis Anterior vs. Tibialis Posterior}]
While both muscles contribute to inversion, \textbf{tibialis anterior} is a primary dorsiflexor and supports medial arch control during swing phase and early stance. In contrast, \textbf{tibialis posterior} acts primarily as a plantarflexor and stabilizer of the medial longitudinal arch during midstance to push-off. Weakness or timing deficits in either may contribute to flatfoot deformity or dynamic valgus.
\end{tcolorbox}

\subsection*{MMT Activities}

\begin{labactivitybox}
\textbf{Activity: Ankle and Foot MMT}
\begin{itemize}
  \item Test plantarflexion, dorsiflexion, inversion, and eversion in gravity-resisted positions.
  \item Use gravity-eliminated positions where needed for weak or painful patients.
  \item Identify peripheral nerve and root level for each movement tested.
\end{itemize}
\end{labactivitybox}

\begin{tcolorbox}[colback=gray!5, colframe=black, title={Clinical Relevance: Inversion, Eversion, and Dynamic Foot Control}]
In open chain, the ankle primarily performs \textbf{inversion} and \textbf{eversion}, controlled by muscles like tibialis posterior/anterior (inverters) and fibularis longus/brevis (evertors).

In closed chain, these motions contribute to \textbf{supination} and \textbf{pronation}—multi-joint composite movements involving the rearfoot, midfoot, and forefoot.

\begin{itemize}
  \item \textbf{Excessive pronation} during stance can result from weak inverters (especially tibialis posterior), leading to medial collapse and altered loading through the tibia, knee, and hip.
  \item \textbf{Insufficient eversion control} reduces lateral ankle stability, increasing risk of inversion sprains during cutting or landing.
\end{itemize}

Evaluating ankle inversion and eversion strength thus provides insight into the body’s ability to control closed-chain supination and pronation. These dynamic control strategies are essential for:
\begin{itemize}
  \item Medial arch support and shock absorption
  \item Frontal plane alignment in single-leg stance and gait
  \item Lateral ankle stability under load
\end{itemize}

This distinction helps bridge isolated strength testing with whole-body movement analysis.
\end{tcolorbox}

\subsection*{DTR: Achilles}

\begin{labactivitybox}
\textbf{Activity: DTR – Achilles Tendon Reflex}

\begin{itemize}
  \item Patient prone with feet over table edge, or seated with legs dangling.
  \item Tap Achilles tendon with foot in slight dorsiflexion; observe plantarflexion.
  \item Use 0–4+ grading scale as above.
  \item Interpret as a screen for tibial nerve and S1 root integrity.
\end{itemize}
\end{labactivitybox}

\section*{Muscle Length Testing}

\begin{labactivitybox}
\textbf{Activity: MLT of Lower Leg and Foot Muscles}
\begin{itemize}
  \item Perform gastrocnemius length test with knee extended, and soleus with knee flexed.
  \item Assess 1st toe extensor flexibility.
  \item Record R1 and R2, compensatory movements, and classify as Normal, Limited, or Excessive.
\end{itemize}
\end{labactivitybox}

\section*{Rehabilitative Ultrasound Imaging of the Ankle}

\begin{labactivitybox}
\textbf{Activity: Achilles Tendon Imaging – RUSI-Based Structural Evaluation\footnote{Based on recommendations from Chimenti et al. (2024), JOSPT Clinical Practice Guideline: Achilles Tendinopathy.}}

\begin{itemize}
  \item \textbf{LAX Views:}
  \begin{itemize}
    \item Scan the Achilles tendon from the myotendinous junction (proximal) through the mid-substance to its insertion on the calcaneus.
    \item Assess tendon continuity, fibrillar architecture, and echogenicity.
    \item Identify pathological features: thickening, hypoechogenic regions, loss of parallel fiber alignment.
  \end{itemize}

  \item \textbf{SAX Views:}
  \begin{itemize}
    \item Image cross-sections at the same levels to confirm tendon contour and symmetry.
    \item Note surrounding anatomical features: Kager’s fat pad, retrocalcaneal bursa, paratenon.
  \end{itemize}

  \item \textbf{Clinical Integration:}
  \begin{itemize}
    \item This activity supports identification of midportion Achilles tendinopathy features, consistent with CPG recommendations.
    \item Use real-time imaging to track changes over time, especially in loading programs.
    \item Save labeled images for structural monitoring: \texttt{Rt/Lt Achilles LAX Mid}, \texttt{Rt/Lt Achilles LAX Ins}, \texttt{Rt/Lt Achilles SAX Mid}.
  \end{itemize}
\end{itemize}
\end{labactivitybox}

\begin{labactivitybox}
\textbf{Activity: Dynamic Ankle RUSI – ATFL/CFL Visualization and Motion Testing}

\textit{Purpose:} To introduce dynamic RUSI techniques for evaluating lateral ankle ligaments during passive movement, aiding in clinical reasoning related to ankle instability and sprain management.

\begin{itemize}
  \item \textbf{Probe Setup:}
  \begin{itemize}
    \item Begin in \textbf{short axis (SAX)} over the anterolateral ankle to identify the anterior talofibular ligament (ATFL).
    \item Slide posteroinferiorly to visualize the calcaneofibular ligament (CFL).
  \end{itemize}
  
  \item \textbf{Dynamic Technique:}
  \begin{itemize}
    \item Apply gentle \textbf{passive plantarflexion and inversion} to stress the ATFL.
    \item For CFL, use \textbf{neutral to dorsiflexed position with inversion}.
    \item Watch for \textbf{ligament excursion}, soft tissue displacement, and any apparent laxity or disruption in fiber continuity.
  \end{itemize}
  
  \item \textbf{Technical Tips:}
  \begin{itemize}
    \item Avoid anisotropy by adjusting probe angle during ligament movement.
    \item Use cine loop capture (if available) to freeze and analyze ligament tensioning through range.
  \end{itemize}

  \item \textbf{Interpretation:}
  \begin{itemize}
    \item Note changes in echotexture or tension across the ligament as motion occurs.
    \item Observe for real-time displacement suggestive of ligamentous insufficiency.
    \item Record probe position, anatomical level, motion type, and dynamic findings.
  \end{itemize}
\end{itemize}

\textit{Clinical Context:} As described in Bush et al.\footnote{Bush et al. (2022). Point of Care Ultrasound Guided Management of Lateral Ankle Sprains. \textit{International Journal of Sports Physical Therapy}}, dynamic RUSI adds value in identifying sprain severity, guiding rehabilitation progression, and improving real-time understanding of ligament mechanics under load.
\end{labactivitybox}

\section*{Pulse Palpation and Imaging}

\begin{labactivitybox}
\textbf{Activity: Distal Arterial Pulse Palpation and Ultrasound}
\begin{itemize}
  \item Palpate dorsalis pedis and posterior tibial arteries.
  \item Use RUSI to confirm location, compressibility, and pulsatility.
\end{itemize}
\end{labactivitybox}

\section*{Clinical Integration}

\begin{itemize}
  \item How did SFMA patterns like SLS, MSF, and ADS help you identify dysfunctions in the foot–ankle–knee chain?
  \item What is the clinical relevance of ankle inversion and eversion strength in controlling dynamic pronation and lateral ankle sprain risk?
  \item How does real-time RUSI of the Achilles tendon contribute to your understanding of tendinopathy staging and rehabilitation progress?
  \item What did you observe during dynamic ultrasound scanning of the lateral ankle ligaments? How might this guide post-injury progression?
  \item How does vascular imaging of the dorsalis pedis or posterior tibial artery confirm your palpation skills and anatomical precision?
\end{itemize}

\section*{Next Steps}

\begin{itemize}
  \item Reinforce your ability to perform MMT for ankle and foot actions, including dorsiflexion, plantarflexion, inversion, and eversion.
  \item Practice ultrasound imaging of the Achilles tendon in LAX and SAX views across its three segments (myotendinous, mid-substance, insertion).
  \item Repeat dynamic ankle scanning to better identify CFL and ATFL motion with passive inversion–plantarflexion.
  \item Review the nerve and root level innervation for each tested muscle to prepare for continued neurovascular reasoning.
  \item Preview Week 8’s lumbopelvic content, focusing on trunk MMT, leg-lowering control, and foundational RUSI of TrA and multifidus.
\end{itemize}

\section*{Checklist Summary}

\begin{itemize}
  \item [$\square$] Interpreted SFMA patterns (ADS, SLS, MSF, MSR) with attention to dorsiflexion, pronation, and medial arch mechanics.
  \item [$\square$] Performed MMT for ankle and foot actions: dorsiflexion, plantarflexion, inversion, eversion, toe flexion.
  \item [$\square$] Identified the peripheral nerve and root level for each tested muscle group.
  \item [$\square$] Performed the Achilles tendon reflex and interpreted findings in relation to S1 and tibial nerve integrity.
  \item [$\square$] Used RUSI to image the Achilles tendon in LAX and SAX at three locations: insertion, mid-substance, myotendinous junction.
  \item [$\square$] Performed introductory dynamic scanning of the lateral ankle to visualize ATFL and CFL during inversion/plantarflexion.
  \item [$\square$] Palpated and imaged the dorsalis pedis or posterior tibial pulse, confirming vessel identity with ultrasound.
\end{itemize}